% preface-ar.tex — Arabic edition of the preface
%
% Original: The Sourdough Framework by Hendrik Kleinwächter,
% the-sourdough-framework.com, CC BY-SA 4.0. The Arabic edition inherits
% CC BY-SA 4.0.
%
% Arabic layer of upstream intro/preface.tex (scaffold structure and the
% \label verbatim; text re-typeset from the recovery reference
% text/01-preface.txt, Phase 28R unit 3 — no retranslation).
% Latin tokens wrapped in \textenglish{} per the record's unwrapped-Latin
% rule (parentheticals carry their parens, per the R1 smoke precedent);
% \qty/\url/math kept from the upstream scaffold. The recovered text's
% البادئة الحامضة (feminine) is re-typeset verbatim per recovery-first;
% the promoted glossary anchors البادئ الحامض (masculine) — flagged in the
% unit report, not silently changed.

\chapter{المقدمة}%
\label{ch:Preface}

إذا كان هناك طعام واحد تُعرف به ألمانيا، فهو على الأرجح الخبز. توجد في
ألمانيا آلاف الأصناف منه، وقد كان صنعُه جزءًا لا يتجزأ من ثقافتنا.

بدأت رحلتي مع الخبز في طفولتي. كانت أمي، وهي أم لثلاثة أبناء، تستغل أيام
السبت دائمًا لخَبز رغيف لذيذ للعائلة. كان خبز توست أبيضًا منفوشًا، وتصنعه
خلال ساعة إلى ساعتين مستخدمةً خميرة جاهزة من المتجر. وبعد أن صرتُ أكثر
خبرة، أدركتُ أنه من الأفضل انتظار قليل قبل تقطيع الخبز، لكننا نحن الأطفال
كنا في تلك الأيام لا نحتمل الانتظار. كانت أمي تقطع لنا بضعة شرائح
مباشرةً من الفرن، وتسرعنا إلى دهن كل شريحة بالزبدة أو المربى. وخلال دقائق
يكون كيلوغرام واحد من الدقيق قد استُهلك. وهكذا صار الخبز جزءًا أساسيًا من
طعامنا الأسبوعي.

كنتُ محظوظًا بأن استطاع والداي تحمّل رحلة تزلج سنوية إلى ألتو أديجي
\textenglish{(Alto Adige)} في شمال إيطاليا. في البلدة الصغيرة المسماة
فالدورا \textenglish{(Valdaora)} كنا نجرّب مطاعم جديدة كل عام، لكننا كنا
ننتهي دائمًا إلى مطعم البيتزا المفضل لدينا. كانت البيتزا هناك مذهلة؛ كان
العجين وحده لذيذاً إلى درجة أننا كنّا نطلب الخبز وحده مع قليل من زيت
الزيتون والملح.

وطبعًا كان سؤالي دائمًا: «أمي، هل يمكننا صنع هذا في البيت من فضلك؟» وهكذا
على مدار السنين صرنا أصدقاء لأصحاب المطعم، وكنا نحصل على إرشادات أكثر
وأكثر حول كيفية صنع عجين البيتزا المثالي. لا توجد فيها مكونات سرية؛ إنها
مجرد دقيق وماء وملح وقليل من الخميرة. كيف يمكن لمزيج بسيط بمقدار هذا أن
يصنع عجين بيتزا لذيذاً إلى هذا الحد؟ وكان والداي من أصحاب العادات
الثابتة، فعادا معنا كل عام، وفي كل عام ازداد اهتمامي. في البيت حاولتُ مع
أمي إعادة صنع الوصفة. جرّبنا الخَبز على حجر وعلى لوح فولاذي. وحاولنا
إضافة زيت إلى العجين وأعشاب إلى صلصة البيتزا. دخلنا في دورة لا تنتهي من
التجارب. ومع ذلك لم نتمكن قط من الاقتراب من التجربة التي عشناها في
الإجازة.

مرّت سنوات، وبدأتُ في النهاية دراستي في مدينة غوتنغن
\textenglish{(Göttingen)} الألمانية الصغيرة. واجهتُ للمرة الأولى مهمة
شراء خبزي بنفسي. ولم يخطر لي أبدًا أن أبدأ الخَبز بنفسي؛ كنت أشتري
رغيفًا جيدًا أثناء تسوقي من السوبرماركت. وكان صنفي المفضل خبزًا يُدعى
\textenglish{\emph{Schwarzbrot: Korn an Korn}}؛ إنه خبز جاودار
\textenglish{(rye)} داكن جدًا وغنيّ، مضافًا إليه توت وبذور عباد الشمس.

وبسذاجة مني، لم أفحص أبدًا من قبل غلافَ ما كنت أشتريه. ثم تغيّر ذلك
يومًا. نظرت إلى الملصق فاندهشت؛ فالخبز الذي بدا صحيًّا كان يتكوّن من
أشياء كثيرة إلى جانب الدقيق والماء. ولم يكن اللون الأسود قادمًا من
الدقيق، بل من سكر مُكرمَل. وذكر الغلاف أنه خبز حامض \textenglish{(sourdough)}،
فلماذا إذن كانت فيه خميرة إضافية؟ ظننتُ أن الخبز الحامض الحقيقي لا يحتاج
إلى خميرة إضافية. وسرعان ما أدركتُ أن في الخبز الذي أشتريه شيئًا ليس على
ما يرام. فتشتُ في خبز السوبرماركت الآخر، فاكتشفتُ أنها جميعًا تحتوي
مكونات لم أسمع بها من قبل. كان ذلك اليوم هو الذي أفقدني الثقة في خبز
السوبرماركت.

في البيت، قررتُ أن أبحث عن الطريقة الصحيحة لصنع الخبز، ويسرّني أنني
تعلّمت أن وصفات صنع البيتزا والخبز متشابهة في الحقيقة إلى حد كبير، مع
وجود فروق أيضًا. فبعض الوصفات مثلًا تطلب خميرة طازجة، وبعضها خميرة جافة.
وحين تعمقتُ في منتديات الإنترنت المختلفة وجميع نقاشاتها العديدة ازددتُ
حيرة. جرّبتُ أنواعًا مختلفة من الدقيق ومارقات مختلفة، من الصنفين العضوي
وغير العضوي معًا. وأدركتُ عندها أنني لا أعرف شيئًا عن صنع الخبز. كانت
الوصفات كثيرًا ما تتعارض فيما بينها فتزيدني حيرة، وبدت أشبه بمجموعة من
خطوات تبدو عشوائية تُتبع فحسب. أما تعليمات الخَبز ودرجات الحرارة فكانت
جميعها مختلفة أيضًا\ldots

وفي هذه الأثناء، وبعد إتمام دراستي، بدأتُ العمل مهندسًا. يواجهنا
المهندسون تحديات كثيرة؛ فالمُصرِّف \textenglish{(compiler)} أو بيئة
التشغيل \textenglish{(runtime)} يصرخان في وجهك بالأخطاء، ومهمتك أن
تكتشف كيف تصلحها. قد يستغرق إصلاح مشكلة بسيطة ساعات، وأحيانًا أيامًا.
وإذا أردت أن تصبح مهندس برمجيات، فعليك أن تنمّي عقلية «لا استسلام»
معينة.

عند كتابة الشيفرة يحتاج مهندسو البرمجيات غالبًا إلى مجموعة من الروتينات
الجاهزة. هذه الروتينات كتبها مهندسون آخرون، ويمكن استخدامها لإيصال
الشيفرة أسرع. ويُعرف هذا الشيفرة المكتوب سلفًا باسم «الإطار»
\textenglish{(framework)}. وفي كثير من الحالات لم يبْنِ هذه الأطر شخص
واحد، بل مهندسون من مختلف أنحاء العالم، كلٌّ منهم يسهم بتحسين الشيفرة
المصدرية وتغييرها. وقد مكّنت الأُطر كثيرًا من الشركات الناجحة. وفي معظم
الحالات تفعل الأُطر ما تدّعيه تمامًا. لكن أحيانًا تواجه مشكلات لا
تفهمها. ففي \qty{99.95}{\percent} من أخطاء البرمجيات يكون المطور هو
المشكلة. لكن أحيانًا يكون في الإطار عيب؛ وعندها على المطور أن يغوص أعمق
ليرى «ماذا» و«لماذا» وراء ما يفعله الإطار. ستحتاج إلى قراءة الشيفرة
المصدرية لمهندسين آخرين، وستُجبَر على فهم «لماذا» تحدث الأمور.

ولكوني غير راضٍ عمّا كنت أخبزه، أخذت عقليتي الهندسية بتّها، واضطررت إلى
الغوص بنفسي في العمق لفهم ما يجري. ولمفاجأتي، لم تُخبرني أي من الوصفات
التي صادفتها «لماذا» ينبغي أن أستخدم مقدارًا $X$ من الماء ومقدارًا $Y$
من الدقيق، أو «لماذا» بالضبط أستخدم الخميرة الطازجة دون الجافة. لماذا
أصفع عجينتي أثناء عجنها \textenglish{(kneading)} على السطح؟ لماذا يكون
الخلاط الكهربائي الثابت \textenglish{(stand mixer)} أفضل من العجن
باليد؟ لماذا أدع العجين يرتاح هذه المدة؟ لماذا يكون توليد البخار
\textenglish{(steaming)} أثناء الخَبز مهمًا؟ وهل أحتاج حقًا إلى شراء قدر
هولندي \textenglish{(dutch oven)} باهظ الثمن لخَبز الخبز؟ ازدادت
المشكلة تعقيدًا حين بدأتُ أقرأ عن البادئة الحامضة
\textenglish{(sourdough starter)}. كان كل ذلك يبدو كسحر أسود. لماذا
صُنعت بعض أنواع البادئ من الفواكه بينما صُنعت أخرى من الدقيق؟ لماذا
تستخدم وصفةٌ القمح بينما تستخدم أخرى الجاودار أو السبلت
\textenglish{(spelt)}؟ كم مرة يجب تغذية البادئ \textenglish{(starter)}؟
كانت الأسئلة التي أحملها وقتئذ تملأ 20 صفحة. كنتُ حائريًا، لكني صرتُ
أكثر تصميمًا على تعلم كيفية صنع خبز لائق في البيت.

ساعدتني ملاحظات الأصدقاء على التحسن مع كل تكرار من خبزي المنزلي. وقارنًا
بالبرمجة، حيث قد تنتظر أحيانًا شهورًا للحصول على هذه الملاحظات، فصنع
الخبز أكثر مباشرةً بكثير. والأفضل أنك يمكنك أكل نجاحاتك (وإخفاقاتك
أيضًا!). ولمفاجأتي، حتى تلك الإخفاقات بدأت تتذوق أفضل من معظم أنواع خبز
المتجر. إن أكل خبز منزلي صنعته بيدك خلال ساعات يمنحك علاقة مختلفة
بطعامك، وكان خَبز الخبز من الصفر بيديّ الخاليتين تغييرًا مرحبًا به بعد
ساعات من العمل على الحاسوب.

واصلتُ التعلم عن عملية التخمير \textenglish{(fermentation)} وتقنيات صنع
الخبز المختلفة. وتعاملتُ مع موضوع البادئة الحامضة بطريقة مشابهة لطريقة
التعامل مع البرمجيات، وبعد سنوات من البحث وتوثيق تقدمي، قررتُ أن حان
الوقت لمشاركة ذلك التقدم مع العالم. عند العمل على مشاريع برمجية من
المهم أن ترى تاريخها وكيف تتغير الشيفرة المصدرية مع الوقت؛ بهذا يمكنك
بسهولة العودة إلى إصدارات سابقة. وكانت هذه الأداة المثالية لتوثيق
وصفاتي، فهي أيضًا تتغير مع كل تكرار جديد. ولمفاجأتي، لقي عملي مفتوح
المصدر عن البادئة الحامضة تقديرَ مهندسين آخرين، وأصبح المشروع مشهورًا
على موقع \textenglish{GitHub}، وهو موقع بُني في الأصل لمشاركة البرمجيات
مفتوحة المصدر.

والآن، عند خَبز خبز رائع تحتاج أيضًا إلى تعلم تقنيات معينة. رأيتُ أنه
سيكون أسهل مشاركتها في هيئة فيديو، وهكذا وُلدت قناتي على
\textenglish{YouTube}. اخترتُ اسم \textenglish{\texttt{The Bread Code}}
ليجسد مقاربتي الهندسية للخبز. استغرق الأمر بعض الوقت حتى أصبت، لكن بعد
أن اخترتُ صورًا مصغرة وعناوين أكثر جاذبية للفيديوهات التي أعدتها، بدأت
القناة تكسب مشاهدين. وأخيرًا، بعد ثلاث سنوات، أخصص يومين من كل أسبوع
لمتابعة شغفي بخَبز الخبز، بينما أواصل في الأيام الثلاثة الأخرى عملي
مهندس برمجيات أكتب شيفرة يومًا بعد يوم.

تملأني أيام الخبز بالفرح والشغف. بالنسبة إليّ لا شيء أفضل من أن أرى كم
شخصًا صنعوا خبزًا مذهلًا بفضل نصائحي وشرحياتي. واصلت المجتمع النمو،
فأنجبت نقاشات وأفكارًا شيقة كثيرة حول موضوع صنع الخبز. هناك دائمًا شيء
جديد نتعلّمه، وأشعر أنه حتى الآن أنا لا أفتح إلا طرف السطح بما أعرفه
وأعلّمه. وهل كنت ستتخيل يومًا أن ذبابات الفاكهة كالنحل، وأنها جزء من
قصة نجاح الخميرة البرية؟ صنعتُ فيديو حاولتُ فيه استنبات أبواغ خميرة
برية قادمة من ذبابات الفاكهة لغرض خَبز الخبز. وقد نجح؛ خرج الخبز مذهل
النتيجة بل وطيّب المذاق! مثل هذه التجارب تشعل اهتمامي الفطري، وإجراؤها
ورؤية الآخرين يتشاركونني الاهتمام يجعلني سعيدًا إلى حد لا يوصف.

وتكمن مشكلة إدارة قناة على \textenglish{YouTube} في أن كل المعلومات
التي تراها تُرشَّح ثم تُقدَّم لك عبر خوارزمية. يقلقني كيف تشكل
الخوارزميات المعلومات الحديثة، لأنها تميل إلى تصنيف المستخدمين في فئات
معينة لا يرون بعدئذ سوى الأخبار المتصلة بذات الفئات الثابتة. وأحد
المقاييس الأساسية المحددة لظهور قناتك هو كم شخصًا نقر على فيديو بعد
عرضه، والمحتوى الذي تصنعه لا يُعرض حتى لكل مشتركي قناتك. وإذا قررت
الخوارزمية أن الفيديو غير جذاب بما يكفي، بدأ محتواك يتلاشى في زوايا
\textenglish{YouTube} المنسية. وحتى لو انتشر الفيديو فوضعًا، ستتوقف
الخوارزمية عن عرضه بمجرد أن تنخفض معدلات التفاعل مع المستخدمين الجدد،
وتتلاشى الفيديوهات القديمة مع الوقت مع ازدياد عامل عقوبة التلاشي. وأنا
أعلم ذلك لأنني طوّرت بنفسي خوارزميات مشابهة كمهندس برمجيات.

ومنذئن قررتُ أن آخذ بعض الوقت بعيدًا عن دورة الخوارزميات لأعمل على شيء
أبعد مدى وأكثر معنى. مهمتي كانت دائمًا أن أشارك معرفتي مع أكبر عدد
ممكن من الناس في العالم؛ ولهذا أيضًا كان محتواي بالإنجليزية لا
بالألمانية. وبعد نقاشات مع أعضاء المجتمع رأيتُ أن كتابة كتاب قد يساعدني
على تحقيق ذلك الهدف. معظم الكتب الموجودة اليوم مجموعات من الوصفات؛ أما
فكرتي فهي أن أقدم لك أساسًا أعمق من المعرفة يمكنك استخدامه لمتابعة
وصفات أخرى. وبمصطلحات البرمجيات، هذا يكون إطار
\textenglish{(bread framework)}.

إن هدفي من هذا الكتاب أن يساعد كل من يواجه مشكلات مع الدقيق أو التخمير
أو الخَبز وغير ذلك. ويجب أن يقدم فهمًا تفصيليًا لسبب كون خطوات معينة
ضرورية، وكيفية تكييفها حين يخطئ الأمر أثناء صنع الخبز. وأرجو أن يكون
هذا العلم في متناول كل إنسان في العالم بغضّ النظر عن ميزانيته، ولذلك لا
أريد أن أتقاضى ثمن الكتاب؛ لهذا قررتُ أن أجعله مفتوح المصدر، وطلبت من
المجتمع دعم عملي بالتبرعات. وملاحظات المجتمع كانت مذهلة حتى الآن، وقد
جمعتُ مالًا أكثر بكثير من المتوقع في البداية. ستبقى النسخة الرقمية من
هذا الكتاب مجانية دائمًا. وهناك أيضًا نسخة غلاف مقوّى متاحة للشراء.
يمكنك قراءة المزيد من التفاصيل هنا:
\url{https://breadco.de/physical-book}

في هذا الكتاب سأحاول أن أكون علميًا قدر الإمكان. لكني لا أدّعي بأي حال
أن الكتاب نفسه سيكون عملًا علميًا. أجريتُ عدة تجارب سأكتب عنها هنا،
لكن لوصف ذلك بأنه علم حقًا ستحتاج على الأرجح إلى تكرار التجربة نفسها
ألف مرة في بيئة مختبرية، وهو ما لم أفعله. سأبذل جهدي مع ذلك لتقديم
مراجع علمية حيث أمكن، وللتمييز بوضوح بين الحقائق والرأي الشخصي.

أرجو أن تستمتع بالقراءة، وأن تتعلم المزيد عن عالم صنع الخبز الساحر،
وأن يمنحك هذا العمل سلسلة الأدوات المتينة التي أتمنى لو كانت بين يديّ
حين بدأتُ رحلتي مع الخبز.

شكرًا لكم.

هندريك
